% ==========================================================================================================
% NOTA: GESTIOME DELLA MODULARITÀ DEL DOCUMENTO E DEI FILE ESTERNI: UTILIZZO DEI COMANDI \input e \include
% ==========================================================================================================
% Nella gestione di progetti complessi, solitamente non è consigliabile generare enormi file di codice sorgente, poiché diventano rapidamente difficili da
% navigare, comprendere e manutenere. Per questo motivo, è buona prassi suddividere il progetto in più file, ossia sfruttare un approccio modulare.
% LaTeX offre due strumenti diversi principali per richiamare file esterni, che hanno scopi e comportamenti molto diversi, ma entrambi fondamentali per
% la creazione di progetti ben strutturati:
%
% 1. COPIA-INCOLLA "TRASPARENTE": \input{percorso/nomefile.tex}
%
% Questo comando agisce come una pura e semplice trascrizione testuale: il contenuto del file esterno viene incollato nel punto esatto della chiamata,
% senza alterare in alcun modo l'impaginazione o il flusso del testo. Grazie a questa proprietà di "trasparenza", \input può essere utilizzato in qualsiasi
% punto del progetto, sia nel preambolo (per caricare pacchetti o impostazioni) sia nel corpo del documento. È inoltre possibile nidificarlo liberamente,
% richiamando un comando \input all'interno di un file già importato.
%
% NOTA: Estensione del file? FACOLTATIVA ma FORTEMENTE CONSIGLIATA. Infatti \input esegue una trascrizione testuale pura, può importare qualsiasi formato
% (es. sorgenti .c, dati .csv). Esplicitare il suffisso .tex rende il codice non ambiguo e professionalmente pulito. Tuttavia può funzionare senza estensione,
% purché il file esista e sia accessibile, e non ci siano omonimie con altri file di diverso formato.
%
% 2.1 GESTIONE MODULARE DI CAPITOLI: \include{percorso/nomefile}
%
% A differenza del primo, \include è progettato per gestire la struttura del corpo del documento. Ad esempio, il suo utilizzo forza un salto di pagina (\clearpage)
% sia prima che dopo l'inserimento del testo (file passato alla sua chiamata) e genera file ausiliari (.aux) dedicati per ottimizzare la gestione dei riferimenti,
% relativamente alla memoria e ai tempi di compilazione. A causa di queste caratteristiche, presenta delle regole di utilizzo molto rigide:
% - Deve essere inserito esclusivamente nel corpo del documento (l'uso nel preambolo genera errori fatali di compilazione).
% - Non supporta alcun tipo di nidificazione.
% - Tipograficamente, è idoneo solo per gli elementi gerarchici di massimo livello, come i capitoli (\chapter) nelle class book (libro), o in quelle basate su
%   di essa, come HKNdocument. L'uso di \include per le sezioni (\section) ad esempio, è altamente sconsigliato, poiché frammetnerebbe molto la lettura e l'estetica
%   del documento, a causa dei grandi spazi vuoti generati dai i salti di pagina che il comando richiama. Proprio per lo stesso motivo, l'uso di \include in class
%   che non supportano i capitoli (\chapter), come article, è altamente sconsigliata.
%
% NOTA: Estensione del file? TASSATIVAMENTE OMESSA. Ingannerebbe il compilatore, che potrebbe cercare di gestire file anomali, come "file.tex.tex", o peggio ancora
% "file.tex.aux" (ricorda che \include genera file ausiliari concatenando l'estensione .aux al nome fornito) compromettendo il funzionamento di \includeonly
% (vedi a seguire). la risoluzione dei riferimenti incrociati.
%
% 2.2 BONUS - STRUMENTO DI COMPILAZIONE SELETTIVA: \includeonly{file1, file2, ...} - UTILE PER TEST E DEBUG
%
% Il principale vantaggio architetturale di \include risiede nella possibilità di utilizzare il comando \includeonly{file1, file2} all'interno del preambolo.
% La stesura di documenti di grandi dimensioni (come tesi o dispense) comporta infatti spesso tempi di compilazione dilatati, specialmente in presenza di grafiche
% vettoriali, circuiti o immagini ad alta risoluzione, tabelle eccetera. Proprio per ovviare a questo ostacolo operativo, LaTeX fornisce l'istruzione
% \includeonly{...}: agisce da filtro sul resto dei comandi \include presenti nel documento, permettendo di compilare rapidamente solo capitoli specifici,
% mantenendo tuttavia intatti e corretti l'indice generale, la numerazione delle pagine e i riferimenti incrociati globali.
% - Posizionamento: Si dichiara esclusivamente all'interno del preambolo.
% - Logica di funzionamento: Il comando \includeonly non inserisce autonomamente alcun file, ma funge esclusivamente da filtro globale di tipo "whitelist"
%   (filtro di sblocco) per i comandi \include già presenti nel corpo del documento, infatti riceve in argomento una lista (separata da virgole) dei file che si desidera compilare
%   (es. \includeonly{capitoli/cap_02, capitoli/cap_05}).
% - Funzionamento tecnico: Il comando impone che il motore TeX, durante la compilazione, vada a comilare esclusivamente i file specificati nella lista, mentre per il resto
%   dei file esclusi, che vada a leggerne  i file ausiliari (.aux) generati nelle compilazioni precedenti.
% - Vantaggio strutturale: A differenza della semplice omissione (tramite %) dei comandi \include nel corpo del testo, \includeonly istruisce il
%   motore TeX su come bypassare la generazione visiva dei capitoli esclusi (le pagine saranno assenti), mantenendo intatti e coerenti numeri di pagina,
%   riferimenti incrociati, indice generale eccetera, come se l'intero documento fosse stato compilato.
% - Prassi di utilizzo: È strettamente uno "strumento da cantiere". Si impiega durante le fasi di redazione per accelerare il rendering del singolo
%   capitolo in lavorazione. A stesura ultimata, la direttiva \includeonly deve essere commentata o rimossa per procedere alla compilazione del
%   volume completo (Release version).
%
% NOTA: Estensione dei file? Come \include.
%
% 0. NOTA GENERALE SUI PERCORSI: È fondamentale utilizzare sempre lo slash (/) come separatore per le directory (es. cartella/file.tex), indipendentemente
% dal sistema operativo in uso. In LaTeX, il backslash (\) è strettamente riservato all'invocazione dei comandi e causerebbe un errore di sintassi.
% ==============================================================================
